\section{Standards and Reference Models}
To guarantee fail-safe operation, railway interlocking systems must conform to internationally recognized safety and engineering standards. These standards not only define technical requirements but also establish the processes, methodologies, and verification practices that ensure signaling systems achieve the highest levels of reliability and integrity.

\subsection{IEC 61508}
IEC 61508 serves as a foundational framework for the functional safety of electrical, electronic, and programmable systems across industries. Its significance lies in introducing the concept of Safety Integrity Levels (SILs), which quantify the probability of system failure and create a structured benchmark for reliability. Beyond classification, the standard emphasizes the importance of lifecycle management, ensuring that safety is considered at every stage of a system’s design, implementation, operation, and eventual decommissioning. For interlocking systems, IEC 61508 acts as the overarching reference, mandating redundancy strategies, fault-tolerant architectures, and rigorous testing protocols that underpin public confidence in railway operations.

\subsection{CENELEC EN 50126/50128/50129}
While IEC 61508 provides a general framework, the CENELEC standards—EN 50126, EN 50128, and EN 50129—translate these principles specifically into the railway domain. EN 50126 outlines the RAMS (Reliability, Availability, Maintainability, and Safety) lifecycle, requiring designers to address not only safety but also the long-term operational resilience of signaling systems. EN 50128 focuses on the software dimension, establishing stringent rules for coding practices, verification, validation, and configuration management. This is particularly relevant for modern computer-based interlocking, where software complexity directly influences safety outcomes. Finally, EN 50129 governs the approval process, bringing together hardware and software assurance to certify the complete system before it enters service. Together, these three standards ensure that railway interlockings are not only technically correct but also systematically verified, traceable, and certifiable in a regulatory environment.

\subsection{UML/SysML Modeling}
In addition to safety standards, contemporary practice increasingly leverages model-based systems engineering (MBSE) as a complementary reference framework. Unified Modeling Language (UML) and Systems Modeling Language (SysML) provide formalized ways to represent interlocking behavior, capturing route logic, state transitions, and failure responses in a structured and visual form. These models allow engineers to validate safety concepts in the early design stages, long before physical implementation, thereby reducing errors and accelerating certification processes. Beyond prototyping, UML and SysML also support traceability by linking safety requirements directly to system functions and test cases, ensuring alignment with the rigorous expectations of standards such as EN 50128. In this sense, modeling serves as a bridge between abstract safety requirements and concrete system design, reinforcing both technical clarity and compliance assurance.
